% =====================================================================
%  Lecture 20 -- 5G Procedures                            (English edition)
% =====================================================================
\documentclass[11pt,a4paper]{article}
\providecommand\ACNroot{../..}
\input{\ACNroot/templates/lecture-en}

\ACNset{number}{20}
\ACNset{title}{5G System Procedures}
\ACNset{shorttitle}{5G procedures}
\ACNset{titleother}{رویه‌های سامانهٔ نسل پنجم}
\ACNset{part}{Part IV -- 5G system and procedures (Lectures 14, 19--22)}
\ACNset{time}{7--9 hours (3.5 h reading, 4--5 h exercises)}
\ACNset{level}{Advanced}
\ACNset{prereq}{Lectures 9, 12, 18 and 19}
\ACNset{gen}{5}

\begin{document}
\ACNcover

% ---------------------------------------------------------------------
\section{Learning objectives}
After studying this lecture you should be able to:
\begin{objectives}
  \item \textbf{Trace} initial registration including 5G-AKA, NAS security and subscription retrieval.
  \item \textbf{Derive} the 5G-AKA verification values and key hierarchy at the level of their inputs and purpose.
  \item \textbf{Trace} PDU session establishment across AMF, SMF, PCF, UPF and the RAN.
  \item \textbf{Explain} UE- and network-triggered service request, RRC\_INACTIVE and Xn/N2 handover.
  \item \textbf{Estimate} procedure latency, AMF signalling load and handover forwarding buffers.
\end{objectives}

% ---------------------------------------------------------------------
\section{Prerequisite knowledge}
This lecture uses the 5GC architecture and identities of Lecture~19, the SBI and NRF of Lecture~18, the EPS-AKA and EPC
procedures of Lectures~9 and~12 (as a comparison) and the PFCP rules of Lecture~13. Procedures are presented in simplified form
following TS~23.502 \cite{ts23502}; optional steps (roaming, interworking, many error cases) are omitted.

\paragraph{Key terms} The vocabulary of this lecture with its Persian equivalents:
\begin{keyterms}
Registration & \fa{ثبت‌نام} & Procedure that makes the UE known to the 5GC (initial, mobility, periodic). \\
5G-AKA & \fa{احراز هویت و توافق کلید ۵G} & Authentication with home-network confirmation of RES*. \\
SEAF & \fa{تابع لنگر امنیتی} & Security anchor in the serving network (co-located with the AMF). \\
NAS SMC & \fa{فرمان حالت امنیتی NAS} & Activation of NAS ciphering and integrity. \\
PDU session establishment & \fa{برقراری نشست PDU} & Creation of a session, its QoS flows and tunnels. \\
N1N2 message transfer & \fa{انتقال پیام N1N2} & AMF service carrying SMF messages to UE (N1) and RAN (N2). \\
Service request & \fa{درخواست خدمت} & Transition from CM-IDLE to CM-CONNECTED. \\
RRC\_INACTIVE & \fa{حالت غیرفعال RRC} & RAN-level suspended state with core context kept. \\
Path switch & \fa{تغییر مسیر} & Moving the N3 tunnel to the target gNB after Xn handover. \\
ABBA & \fa{پارامتر ضدحمله به پایین‌برداری} & Anti-bidding-down parameter bound into K\textsubscript{AMF}. \\
\end{keyterms}

% ---------------------------------------------------------------------
\section{Introduction and motivation}
An architecture becomes concrete only through its procedures: the ordered exchanges that make a UE known to the network, secure
its signalling, give it a data session and keep that session alive while the UE sleeps and moves. In the 5GC these procedures
combine three kinds of signalling: NAS between UE and AMF (TS~24.501 \cite{ts24501}), NGAP between RAN and AMF (TS~38.413
\cite{ts38413}) and service operations between NFs (Lecture~18), plus PFCP towards the UPF.

Understanding the procedures is essential for troubleshooting (where did an attach fail?), for dimensioning (how many messages
does each NF handle?) and for latency engineering (how long until the first packet flows?). This lecture traces the four central
procedures and quantifies their cost.

% ---------------------------------------------------------------------
\section{Fundamental concepts}

\subsection{Types of registration}
\begin{itemize}
  \item \textbf{Initial registration}: the UE is RM-DEREGISTERED and registers for the first time (or after a long absence).
  \item \textbf{Mobility registration update}: the UE enters a tracking area outside its registration area, or its capabilities or
        requested slices change.
  \item \textbf{Periodic registration update}: on expiry of the periodic timer, to confirm reachability.
  \item \textbf{Emergency registration}: limited service for emergency calls.
\end{itemize}

\subsection{Security anchors}
5G-AKA \cite{ts33501} differs from EPS-AKA in one important respect: the \emph{home network} confirms the authentication. The serving
network's SEAF checks a hashed expected response (HXRES*) for a first verification, but the AUSF in the home network verifies RES*
itself before releasing the anchor key $K_{\mathrm{SEAF}}$ and the SUPI. A visited network therefore cannot claim that a subscriber was
present without proof -- an answer to fraud concerns in roaming.

\subsection{Session establishment roles}
The AMF only relays session signalling: it selects an SMF and transports N1 (NAS SM) and N2 (NGAP resource set-up) messages for it. The
SMF owns the session: it obtains policy from the PCF, selects and programs the UPF over N4, and asks the AMF to deliver the NAS accept to
the UE and the resource request to the RAN in a single \emph{N1N2 message transfer}.

% ---------------------------------------------------------------------
\section{Detailed technical explanation}

\subsection{Initial registration}
The UE sends a Registration Request with its SUCI (or a 5G-GUTI if it has one), requested slices (NSSAI) and capabilities
(\cref{fig:reg}) \cite{ts23502}. The RAN selects an AMF (from the GUAMI in the 5G-S-TMSI or by load) and forwards the message in an NGAP
Initial UE Message. If the UE gave a 5G-GUTI from another AMF, the new AMF fetches the context from the old AMF
(Namf\_Communication\_UEContextTransfer). The AMF then authenticates the UE (5G-AKA, \cref{sec:aka}), activates NAS security, registers
with the UDM (Nudm\_UECM\_Registration), retrieves the access and mobility subscription (Nudm\_SDM\_Get) and subscribes to changes, creates
an AM policy association with the PCF, and, if needed, asks the NSSF to select slices. Registration Accept carries the new 5G-GUTI, the
registration area and the allowed NSSAI; the UE acknowledges with Registration Complete.

\begin{acnfigure}
  \resizebox{\linewidth}{!}{\acnfig{lecture-20/registration}}
  \caption{Initial registration with 5G-AKA (simplified) \cite{ts23502,ts33501}.}
  \label{fig:reg}
\end{acnfigure}

\subsection{5G-AKA}\label{sec:aka}
The UDM generates an authentication vector with the same functions $f_1$--$f_5$ as in Lecture~11, but derives 5G-specific values bound to
the \emph{serving network name} (SNN, of the form \nolinkurl{5G:mnc<MNC>.mcc<MCC>.3gppnetwork.org}) \cite{ts33501}:
\begin{subequations}\label{eq:aka}
\begin{align}
  \mathrm{XRES}^{*} &= \KDF\bigl(\mathrm{CK}\concat\mathrm{IK},\ \mathrm{SNN},\ \mathrm{RAND},\ \mathrm{XRES}\bigr)\ \text{(128 bits)}, \label{eq:xres}\\
  \mathrm{HXRES}^{*} &= \mathrm{SHA\text{-}256}\bigl(\mathrm{RAND}\concat\mathrm{XRES}^{*}\bigr)\ \text{(truncated to 128 bits)}, \label{eq:hxres}\\
  K_{\mathrm{AUSF}} &= \KDF\bigl(\mathrm{CK}\concat\mathrm{IK},\ \mathrm{SNN},\ \mathrm{SQN}\xor\mathrm{AK}\bigr), \label{eq:kausf}
\end{align}
\end{subequations}
followed by
\[
\begin{aligned}
  K_{\mathrm{SEAF}} &= \KDF(K_{\mathrm{AUSF}}, \mathrm{SNN}), \qquad
  K_{\mathrm{AMF}} = \KDF(K_{\mathrm{SEAF}}, \mathrm{SUPI}, \mathrm{ABBA}),\\
  K_{\mathrm{gNB}} &= \KDF(K_{\mathrm{AMF}}, \text{UL NAS COUNT}, \ldots).
\end{aligned}
\]
The AUSF stores XRES* and gives the SEAF only HXRES*: the SEAF can check whether
the hash of RAND and RES* computed as in \cref{eq:hxres} matches, but cannot compute RES* itself, so it must forward RES* to the
home network for the final decision.

\subsection{PDU session establishment}
The UE sends a PDU Session Establishment Request (with a new PDU session ID, DNN, S-NSSAI, session type and SSC mode) inside an UL NAS
Transport message (\cref{fig:pdu}). The AMF selects an SMF (NRF discovery by DNN and slice) and calls Nsmf\_PDUSession\_CreateSMContext.
The SMF retrieves session-management subscription data from the UDM, creates an SM policy association with the PCF (which returns PCC rules),
selects a UPF and establishes the N4 session (PDRs, FARs, QERs, URRs), obtaining the uplink tunnel endpoint. It then calls
Namf\_Communication\_N1N2MessageTransfer with the NAS accept (QoS rules, IP address) for the UE and the N2 resource set-up information (QoS
flows and profiles, UPF tunnel endpoint) for the gNB. The gNB configures data radio bearers with RRC Reconfiguration and returns its
downlink tunnel endpoint, which the SMF installs in the UPF with an N4 modification. QoS flows are now active in both directions.

\begin{acnfigure}
  \resizebox{\linewidth}{!}{\acnfig{lecture-20/pdu-session}}
  \caption{PDU session establishment (simplified) \cite{ts23502}.}
  \label{fig:pdu}
\end{acnfigure}

\subsection{Service request and paging}
A UE in CM-IDLE with uplink data sends a Service Request (with its 5G-S-TMSI) in RRC set-up; the AMF re-activates the user plane of the
selected PDU sessions through the SMFs (UpdateSMContext, N4 modification) and the RAN (N2 Initial Context Setup). For downlink data, the UPF
buffers (or the SMF buffers, by FAR configuration) and notifies the SMF over N4; the SMF calls N1N2MessageTransfer; the AMF pages the UE in
its registration area; the UE answers with a Service Request. A UE in RRC\_INACTIVE skips most of this: the core context and N3 tunnels are
kept, the RAN pages within its notification area and resumes the RRC connection locally, so no core signalling is needed unless the UE has
moved out of the area.

\subsection{Handover}
With an Xn interface between gNBs, the source decides, prepares the target over Xn, commands the UE, forwards buffered data; the target then
sends an N2 Path Switch Request so that the AMF and SMF move the N3 tunnel to the target, with end markers on the old path (\cref{fig:xn}).
Without Xn (or when the AMF or UPF must change), an \emph{N2 handover} is prepared through the AMF(s), which also handles AMF relocation
and inter-system handover to LTE.

\begin{acnfigure}
  \resizebox{\linewidth}{!}{\acnfig{lecture-20/xn-handover}}
  \caption{Xn handover with path switch (simplified).}
  \label{fig:xn}
\end{acnfigure}

% ---------------------------------------------------------------------
\section{Mathematical models}

\subsection{Procedure latency}\label{sec:lat}
\paragraph{Model} A procedure is a chain of sequential exchanges; parallel branches count only once. With $n_{\mathrm{N}}$ UE--AMF round trips of
duration $T_{\mathrm{N}}$ (radio and RAN included) and $n_{\mathrm{S}}$ sequential core exchanges of duration $T_{\mathrm{S}}$,
\begin{equation}\label{eq:latency}
  T_{\mathrm{proc}} \approx n_{\mathrm{N}}\,T_{\mathrm{N}} + n_{\mathrm{S}}\,T_{\mathrm{S}} .
\end{equation}
\paragraph{Example} Initial registration (\cref{fig:reg}) has about 6 UE--AMF round trips (registration, authentication, security mode, accept
plus RRC) and 8 sequential SBI exchanges. With $T_{\mathrm{N}} = 15$\,ms and $T_{\mathrm{S}} = 3$\,ms: $T_{\mathrm{proc}} \approx 90 + 24 = 114$\,ms. The radio
round trips dominate; core optimisation helps mainly when NFs are far apart or overloaded.

\subsection{Signalling load at the AMF}\label{sec:load}
\paragraph{Model} If procedure type $k$ occurs $r_k$ times per UE per hour and costs $m_k$ messages at the AMF (NAS, NGAP and SBI counted), an AMF
serving $N$ UEs handles
\begin{equation}\label{eq:load}
  \Lambda = \frac{N}{3600}\sum_k r_k\, m_k \quad \text{messages/s}.
\end{equation}
\paragraph{Example} Illustrative values for $N = 10^6$: registrations 0.5/h at 14 messages, service requests 20/h at 6, Xn handovers 2/h at 4, PDU
session establishments 0.2/h at 10: $\Lambda = 10^6\times137/3600 = 3.8\times10^4$ messages/s, of which service requests make up 88\,\%
(\cref{fig:load}). If 70\,\% of idle-to-active transitions use RRC\_INACTIVE (no core messages), $\Lambda$ falls to $1.5\times10^4$ messages/s.
\paragraph{Interpretation} Frequent short sessions -- smartphones and IoT devices waking for a few packets -- dominate core signalling; RRC\_INACTIVE
is primarily a signalling-reduction feature.

\begin{acnfigure}
  \acnfig{lecture-20/amf-load}
  \caption{AMF message load per procedure type for the illustrative model of \cref{sec:load}.}
  \label{fig:load}
\end{acnfigure}

\subsection{Forwarding buffer during handover}\label{sec:fwd}
\paragraph{Model} Between the moment the source stops transmitting and the moment the path switch is complete, downlink data at rate $R$ are
forwarded over Xn and buffered at the target until the UE is ready. With interruption $T_{\mathrm{s}}$,
\begin{equation}\label{eq:fwd}
  B = R\,T_{\mathrm{s}} .
\end{equation}
\paragraph{Example} $R = 100$\,Mbit/s, $T_{\mathrm{s}} = 20$\,ms: $B = 2$\,Mbit $= 250$\,kB per UE. For an eMBB cell with several fast users the target
needs megabytes of buffer; for URLLC, dual active protocol stack handover avoids the interruption altogether.

\subsection{Authentication security}\label{sec:sec}
\paragraph{Model} An attacker without $K$ who guesses RES* (128 bits) succeeds with probability $2^{-128}$ per attempt; one who has observed HXRES*
must invert a 128-bit truncated hash. The expected number of attempts for a brute-force search over $b$ bits is
\begin{equation}\label{eq:brute}
  \E[A] = 2^{b-1}.
\end{equation}
\paragraph{Example} $b = 128$: $1.7\times10^{38}$ attempts -- infeasible; the security of 5G-AKA rests on the secrecy of $K$ and the freshness of SQN,
not on the length of the verification values.

% ---------------------------------------------------------------------
\section{Network architecture and diagrams}
The figures of this lecture are \emph{message sequence charts}: vertical lifelines for UE, RAN and NFs, horizontal arrows for messages.
They abstract from transport (SCTP for NGAP, HTTP/2 for SBI, UDP for PFCP), from intermediate SCPs and from parallelism. In a real trace the
same procedure appears as dozens of packets on several interfaces; reading traces efficiently means mapping them back to these charts.

% ---------------------------------------------------------------------
\section{Step-by-step procedures}

\begin{procedurebox}{Initial registration}\label{proc:reg}
\begin{steps}
  \item UE $\to$ AMF: Registration Request (SUCI, requested NSSAI, capabilities) via RRC set-up and NGAP Initial UE Message.
  \item AMF $\to$ AUSF $\to$ UDM: authentication request; the UDM de-conceals the SUCI and returns a 5G home-environment vector.
  \item AMF (SEAF) $\to$ UE: Authentication Request; UE verifies AUTN and returns RES*; SEAF checks HXRES*; AUSF checks RES* and returns
        SUPI and $K_{\mathrm{SEAF}}$.
  \item NAS Security Mode Command/Complete; the AMF derives $K_{\mathrm{AMF}}$ and NAS keys.
  \item AMF registers with the UDM and retrieves subscription data; creates the AM policy association; selects slices.
  \item Registration Accept (5G-GUTI, registration area, allowed NSSAI); Registration Complete.
\end{steps}
\end{procedurebox}

\begin{procedurebox}{PDU session establishment}\label{proc:pdu}
\begin{steps}
  \item UE $\to$ AMF: PDU Session Establishment Request (DNN, S-NSSAI, type, SSC mode) in UL NAS Transport.
  \item AMF selects an SMF and calls CreateSMContext; SMF fetches SM subscription, creates the SM policy association.
  \item SMF selects a UPF and sends N4 Session Establishment with the rules; receives the UL tunnel endpoint.
  \item SMF $\to$ AMF: N1N2MessageTransfer (N1: Accept with QoS rules and IP address; N2: QoS flows, UL TEID).
  \item AMF $\to$ gNB: PDU Session Resource Setup; gNB configures DRBs and returns the DL TEID.
  \item AMF $\to$ SMF: UpdateSMContext; SMF $\to$ UPF: N4 modification with the DL TEID.
\end{steps}
\end{procedurebox}

% ---------------------------------------------------------------------
\section{Worked examples}

\begin{examplebox}{Who checks what in 5G-AKA}\label{ex:aka}
\textbf{Problem.} A visited network's AMF receives RES* from the UE. Explain which checks it can perform and why it needs the AUSF.

\textbf{Solution.} The SEAF has HXRES* and RAND; it computes $\mathrm{SHA\text{-}256}(\mathrm{RAND}\concat\mathrm{RES}^*)$ (truncated) and compares with
HXRES* (\cref{eq:hxres}) -- this filters wrong responses locally. It cannot derive XRES* (it lacks CK, IK) and does not receive the SUPI
or $K_{\mathrm{SEAF}}$ until the AUSF has compared RES* with XRES*. The home network thus confirms the authentication.
\end{examplebox}

\begin{examplebox}{Time to first packet}\label{ex:first}
\textbf{Problem.} After power-on, a UE registers and then establishes a PDU session. The PDU session has 3 UE--AMF round trips (request, RRC
reconfiguration with accept) and 6 sequential SBI/N4 exchanges. With the values of \cref{sec:lat}, when can the first packet flow?

\textbf{Solution.} Registration 114\,ms; PDU session $3\times15 + 6\times3 = 63$\,ms; total about 180\,ms (plus cell search and system information
acquisition, typically longer). Running the PDU session request piggybacked on registration (allowed for some cases) or pre-establishing
sessions shortens it.
\end{examplebox}

\begin{examplebox}{Dimensioning with RRC\_INACTIVE}\label{ex:inactive}
\textbf{Problem.} An AMF instance handles $2\times10^4$ messages/s. How many instances are needed for the load of \cref{sec:load} with and without
RRC\_INACTIVE, at 70\,\% target utilisation?

\textbf{Solution.} Without: $3.8\times10^4/(0.7\times2\times10^4) = 2.7 \to 3$ instances; with: $1.47\times10^4/1.4\times10^4 = 1.05 \to 2$ instances (plus
redundancy in both cases).
\end{examplebox}

% ---------------------------------------------------------------------
\section{Technology comparison}
\begin{table}[htbp]
\centering\small
\caption{EPS and 5GS procedures compared.}\label{tab:cmp}
\rowcolors{2}{acnmist}{white}
\begin{tabularx}{\linewidth}{@{}l Y Y@{}}
\toprule
\rowcolor{white} & EPS & 5GS \\
\midrule
Attach / registration & Attach always creates a default bearer & Registration without session; PDU sessions separately \\
Authentication        & EPS-AKA, decision in serving network & 5G-AKA or EAP-AKA$'$, home confirmation \\
Identity              & IMSI may be sent in clear          & SUCI (concealed) \\
Session set-up        & MME $\to$ S-GW $\to$ P-GW (GTPv2-C)  & AMF relays; SMF drives PCF, UPF (SBI, PFCP) \\
Idle optimisation     & ECM-IDLE only                      & CM-IDLE and RRC\_INACTIVE \\
Handover path switch  & MME $\to$ S-GW (S11)               & AMF $\to$ SMF $\to$ UPF (SBI, N4) \\
\bottomrule
\end{tabularx}
\end{table}

\paragraph{Always-on versus on-demand} The EPS attach always created a default bearer (``always-on IP''); 5G separates registration from
session establishment. A device can register (e.g.\ for SMS or for a later session) without consuming user-plane resources, and sessions are
created per DNN and slice as needed.

% ---------------------------------------------------------------------
\section{Key takeaways}
\begin{takeawaybox}
\begin{itemize}
  \item Registration = identification (SUCI/5G-GUTI) + 5G-AKA + NAS SMC + UDM registration and subscription + policy + accept.
  \item 5G-AKA binds keys to the serving network name; SEAF checks HXRES*, the home AUSF confirms RES*.
  \item PDU session: AMF relays, SMF orchestrates PCF and UPF, N1N2 transfer delivers NAS accept and N2 set-up together.
  \item Service requests dominate AMF load; RRC\_INACTIVE removes most of them from the core.
  \item Latency $\approx n_{\mathrm{N}}T_{\mathrm{N}} + n_{\mathrm{S}}T_{\mathrm{S}}$; load $= (N/3600)\sum r_k m_k$; forwarding buffer $= RT_{\mathrm{s}}$.
\end{itemize}
\end{takeawaybox}

% ---------------------------------------------------------------------
\section{Exercises}

\begin{exercisebox}{1}{Order of steps}\label{exr:order}
Why must NAS security be activated before the AMF sends the Registration Accept? What could an attacker do otherwise?
\end{exercisebox}

\begin{exercisebox}{1}{N1N2 message transfer}\label{exr:n1n2}
Explain why the SMF sends both the NAS accept and the N2 resource information to the AMF in one service operation rather than in two.
\end{exercisebox}

\begin{exercisebox}{2}{Latency}\label{exr:latency}
A mobility registration update needs 2 UE--AMF round trips and 4 sequential SBI exchanges (context transfer from the old AMF, UDM update,
policy update, subscription retrieval). Compute its latency with $T_{\mathrm{N}} = 15$\,ms and $T_{\mathrm{S}} = 3$\,ms, and with $T_{\mathrm{S}} = 12$\,ms if the
old AMF is in a distant data centre.
\end{exercisebox}

\begin{exercisebox}{2}{Forwarding buffer}\label{exr:fwd}
A cell has 5 users each receiving 300\,Mbit/s when a train carries them all into a neighbouring cell simultaneously; the interruption is 30\,ms.
How much buffer does the target gNB need? What would reduce it?
\end{exercisebox}

\begin{exercisebox}{3}{IoT signalling}\label{exr:iot}
$10^6$ IoT devices per AMF send one small report every 15 minutes. Each report needs a service request (6 AMF messages) if the device returns to
CM-IDLE, or none if it stays in RRC\_INACTIVE (resumed in the RAN). Compute the AMF load in both cases, and discuss what RRC\_INACTIVE costs the
RAN instead.
\end{exercisebox}

\begin{exercisebox}{3}{Authentication in roaming}\label{exr:roam}
A visited network claims to have served a roaming subscriber for 2 hours and bills the home operator. With EPS-AKA and with 5G-AKA, what evidence
does the home network have that the subscriber was really authenticated in the visited network? Relate your answer to \cref{eq:aka}.
\end{exercisebox}

\subsection*{Solutions}
\addcontentsline{toc}{subsection}{Solutions}

\begin{solutionbox}{\ref{exr:order}}
The Registration Accept carries the new 5G-GUTI, the allowed slices and the registration area. Without integrity protection an attacker could
modify these (e.g.\ restrict the UE's services or force it to reveal itself), and without ciphering the GUTI--SUCI link would be visible,
undermining identity privacy. The security mode procedure also confirms the UE's security capabilities, preventing bidding-down attacks.
\end{solutionbox}

\begin{solutionbox}{\ref{exr:n1n2}}
The NAS accept and the RAN resource set-up belong together: the RAN must create radio bearers for the QoS flows the UE is told about, and the AMF
can deliver both in one NGAP message (PDU Session Resource Setup with the NAS PDU inside). One operation keeps them consistent and saves a round
trip; if the RAN rejects the resources, the AMF can report the failure to the SMF without the UE ever seeing an inconsistent accept.
\end{solutionbox}

\begin{solutionbox}{\ref{exr:latency}}
$2\times15 + 4\times3 = 42$\,ms; with $T_{\mathrm{S}} = 12$\,ms: $30 + 48 = 78$\,ms. Distant core NFs can dominate procedures with many sequential
exchanges -- an argument for locality-aware NF selection.
\end{solutionbox}

\begin{solutionbox}{\ref{exr:fwd}}
$B = 5\times300\times10^6\times0.03/8 = 5.6$\,MB. Reductions: shorter interruption (RACH-less or conditional handover, prepared target), dual active
protocol stack (no interruption), or reducing the forwarded volume by letting TCP slow down; on trains, a moving relay or group handover helps.
\end{solutionbox}

\begin{solutionbox}{\ref{exr:iot}}
Reports: $10^6/900 = 1111$\,s$^{-1}$. CM-IDLE: $1111\times6 = 6667$ AMF messages/s; RRC\_INACTIVE: none for the reports (only periodic RAN notification
area updates, also handled in the RAN). The RAN must keep the UE contexts and security state for all inactive UEs (memory) and handle RAN paging
and resume; when a UE moves out of the RAN notification area, a RAN-based update is needed.
\end{solutionbox}

\begin{solutionbox}{\ref{exr:roam}}
With EPS-AKA the serving network compares RES with XRES itself; the home network only knows that it sent authentication vectors -- it cannot prove
the UE responded. With 5G-AKA the AUSF receives RES* and compares it with XRES* (\cref{eq:xres}); RES* can only be computed by a USIM holding $K$ in
response to that RAND and bound to that serving network name. The home network thus has cryptographic evidence that the subscriber was authenticated
in that network at that time.
\end{solutionbox}

% ---------------------------------------------------------------------
\section{Further reading}
The procedures are specified in TS~23.502 \cite{ts23502}, with NAS in TS~24.501 \cite{ts24501}, NGAP in TS~38.413 \cite{ts38413} and security
(5G-AKA, key hierarchy, SUCI) in TS~33.501 \cite{ts33501}; the SBI operations used are in TS~29.500 \cite{ts29500} and the NRF in TS~29.510
\cite{ts29510}. RRC\_INACTIVE and handover on the RAN side are in TS~38.300 \cite{ts38300}. Rommer et al.\ \cite{rommer2020core} walk through the
procedures in detail.

\printbibliography[heading=subbibliography,title={\ACNstrRefs}]
\end{document}
